\documentclass{article}

% ICML 2025 style
\usepackage[accepted]{icml2025}

% Standard packages
\usepackage[utf8]{inputenc}
\usepackage[T1]{fontenc}
\usepackage{hyperref}
\usepackage{url}
\usepackage{booktabs}
\usepackage{amsmath}
\usepackage{amssymb}
\usepackage{graphicx}
\usepackage{microtype}
\usepackage{xcolor}
\usepackage{multirow}
\usepackage{enumitem}

\graphicspath{{./}}

\tolerance=800
\emergencystretch=12pt

% -----------------------------------------------------------------------
\icmltitlerunning{Cross-Lingual Hate Speech Detection: English to Roman Urdu}

\begin{document}

\twocolumn[
\icmltitle{Cross-Lingual Hate Speech Detection:\\
           English to Roman Urdu Transfer Learning}

\begin{icmlauthorlist}
\icmlauthor{Muteeullah Baig}{uu}
\icmlauthor{Rana Muhammad Hamza}{uu}
\end{icmlauthorlist}

\icmlaffiliation{uu}{AI-600 Deep Learning, LUMS, Spring 2026}
\icmlcorrespondingauthor{Muteeullah Baig}{25280036@lums.edu.pk}
\icmlcorrespondingauthor{Rana Muhammad Hamza}{25280103@lums.edu.pk}

\icmlkeywords{hate speech detection, cross-lingual transfer, multilingual transformers,
              Roman Urdu, XLM-RoBERTa, few-shot learning, prompt-based fine-tuning, fairness}

\vskip 0.3in
]

\printAffiliationsAndNotice{}

% Code: \url{https://github.com/MuteeullahBaig/AI600-HateSpeech-CrossLingual}

% -----------------------------------------------------------------------
\begin{abstract}
We study cross-lingual transfer of hate speech detection from English to Roman
Urdu, a severely under-resourced informal script used on Pakistani social media.
Fine-tuning three multilingual transformers (mBERT, XLM-R-base, XLM-R-large)
on English data yields strong English performance (XLM-R-large: hate~F1~$=$~0.509)
but near-zero zero-shot transfer to Roman Urdu (hate~F1~$<$~0.07 for all models).
Standard few-shot fine-tuning ($K \le 64$) collapses to all-hate predictions due to
class distribution mismatch between source and target languages.
Prompt-based fine-tuning with an MLM verbalizer substantially mitigates this
collapse: macro~F1 improves from 0.349 to 0.745 at $K=64$ vs.\ standard fine-tuning.
Supervised TF-IDF$+$FuzzyWuzzy baselines and full XLM-R-large Urdu supervision
reach hate~F1~$=$~0.857 and 0.895 respectively, bounding the remaining cross-lingual
gap. Fairness analysis reveals a 5$\times$ higher false positive rate on
religion-mentioning text.
\end{abstract}

% -----------------------------------------------------------------------
\section{Introduction}
\label{sec:intro}

Hate speech on social media disproportionately affects non-English-speaking
communities, yet content moderation systems are largely built for English.
Roman Urdu --- Urdu written in the Latin alphabet --- is the predominant informal
script on Pakistani social media and presents unique challenges: mixed
English--Urdu vocabulary, non-standard orthography, and culturally specific
hate speech patterns absent from English training corpora.

This work investigates \textit{cross-lingual transfer learning} for hate speech
detection, treating English as the source language and Roman Urdu as the target.
We address four research questions:

\begin{itemize}[leftmargin=*,nosep]
  \item \textbf{RQ1 (Backbone):} How does model capacity affect English detection
        and cross-lingual transfer?
  \item \textbf{RQ2 (Zero-shot):} Can a model trained on English detect Roman Urdu
        hate speech without target-language supervision?
  \item \textbf{RQ3 (Few-shot):} How do $K$-shot Urdu examples affect transfer,
        and can prompt-based fine-tuning mitigate degenerate collapse?
  \item \textbf{RQ4 (Fairness):} Does the transferred model exhibit differential
        false positive rates across demographic groups?
\end{itemize}

% -----------------------------------------------------------------------
\section{Related Work}
\label{sec:related}

\paragraph{Hate speech detection.}
\citet{davidson2017automated} introduced an influential English Twitter dataset and
demonstrated that distinguishing hate speech from offensive language is challenging
even for humans due to inter-annotator disagreement on borderline cases.
\citet{fortuna2018survey} survey automatic methods and highlight that most work
focuses on English data, with performance often inflated by dataset-specific
lexical cues rather than generalizable linguistic understanding.

\paragraph{Multilingual transformers.}
BERT~\citep{devlin2019bert} introduced transformer pre-training.
Its multilingual variant (mBERT) demonstrated surprising zero-shot cross-lingual
transfer~\citep{pires2019multilingual}.
XLM-RoBERTa (XLM-R)~\citep{conneau2020unsupervised} trained on 2.5\,TB of
Common Crawl data across 100 languages, consistently outperforming mBERT on
cross-lingual benchmarks.

\paragraph{Cross-lingual hate speech.}
\citet{nozza2021exposing} found that zero-shot cross-lingual hate speech transfer
degrades substantially for distant language pairs.
\citet{ousidhoum2019multilingual} analysed multilingual hate speech across European
languages.
\citet{jiang2024crosslingual} survey 82 multilingual datasets across 10 language
families and identify cultural nuance and limited labelled data as the two persistent
bottlenecks in cross-lingual transfer.

\paragraph{Roman Urdu hate speech.}
\citet{rizwan2020hate} introduced the first Roman Urdu hate speech dataset and
showed that transfer from English embeddings outperforms training from scratch.
\citet{ullah2024fuzzy} showed that fusing TF-IDF features with FuzzyWuzzy
string-matching scores outperforms weak multilingual transformers on Roman Urdu
cybercrime classification using classical ML classifiers.
Neither work studies transfer from a strong multilingual model such as XLM-R-large
in the zero- or few-shot setting.

\paragraph{Prompt-based fine-tuning.}
\citet{ullah2025prompt} replace the standard classification head with a masked
language modelling (MLM) objective, wrapping each input in a template and mapping
the predicted token to a class label via a verbalizer.
With only 32 examples per class, prefix-prompt XLM-R matches standard fine-tuning
on 70--80\% of data across eight high-resource languages.
\citet{maab2026prompt} extend prompting to offensive Urdu language detection using
LLMs.
We adapt this approach to the harder cross-lingual setting to test whether the
MLM-head advantage mitigates the few-shot collapse we observe.

% -----------------------------------------------------------------------
\section{Dataset}
\label{sec:data}

\paragraph{English source data.}
We use the Davidson et al.\ (2017) dataset from HuggingFace
(\path{tdavidson/hate_speech_offensive}), binarised into hate (class~0) and
non-hate (classes~1--2). Table~\ref{tab:data} reports statistics.

\paragraph{Roman Urdu target data.}
We use \path{community-datasets/roman_urdu_hate_speech} from HuggingFace,
containing Twitter posts annotated by native speakers.

\begin{table}[h]
\centering
\caption{Dataset statistics.}
\label{tab:data}
\vskip 0.05in
\begin{small}
\begin{tabular}{llrcc}
\toprule
Split & Lang & $n$ & Non-hate & Hate \\
\midrule
Train & EN   & 21,159 & 94.8\% & 5.2\%  \\
Dev   & EN   &  2,612 & 94.7\% & 5.3\%  \\
Train & UR   &  7,208 & 46.6\% & 53.4\% \\
Dev   & UR   &    800 & 46.5\% & 53.5\% \\
\bottomrule
\end{tabular}
\end{small}
\end{table}

The English data is severely imbalanced (5.2\% hate), while the Roman Urdu corpus
is nearly balanced (53.5\% hate). This distribution mismatch is central to
understanding the few-shot transfer failure (Section~\ref{sec:results}).

% -----------------------------------------------------------------------
\section{Methodology}
\label{sec:method}

\subsection{Model Architectures}

We compare three multilingual transformers:
\textbf{mBERT}~\citep{devlin2019bert} (179\,M params, 12 layers),
\textbf{XLM-R-base}~\citep{conneau2020unsupervised} (278\,M params, 12 layers), and
\textbf{XLM-R-large} (560\,M params, 24 layers).
All are fine-tuned end-to-end by appending a two-class linear head on the
\texttt{[CLS]} representation.

\subsection{English Fine-Tuning Protocol}

All three models are trained on the English dataset using the HuggingFace Trainer
API with class-weighted cross-entropy (Table~\ref{tab:hparams}).
The English training set is 94.8\% non-hate; without weighting, models achieve
near-zero hate~F1 by predicting all non-hate.
We upweight the hate class by~8 and select the best epoch by dev-set hate~F1.

\begin{table}[h]
\centering
\caption{English training hyperparameters.}
\label{tab:hparams}
\vskip 0.05in
\begin{small}
\begin{tabular}{lccc}
\toprule
Hyperparameter & mBERT & XLM-R-b & XLM-R-l \\
\midrule
Learning rate     & 2e-5  & 2e-5  & 1e-5  \\
Batch size        & 32    & 32    & 32    \\
Epochs            & 5     & 5     & 5     \\
fp16              & Yes   & Yes   & Yes   \\
Grad.\ ckpt.      & No    & No    & Yes   \\
Class weights     & \multicolumn{3}{c}{[1.0, 8.0]} \\
\bottomrule
\end{tabular}
\end{small}
\end{table}

\subsection{Classical ML Baselines}

Following \citet{ullah2024fuzzy}, we implement TF-IDF (unigrams+bigrams,
top 50{,}000 features) combined with Logistic Regression (LR) or Random Forest
(RF), with and without FuzzyWuzzy string-matching features.
FuzzyWuzzy computes token set ratio and partial ratio scores between each tweet
and a lexicon of known hate speech terms, providing a soft lexical similarity
signal complementary to TF-IDF bag-of-words.
All classifiers use \texttt{class\_weight=`balanced'} and are trained on the
full Roman Urdu training set (7{,}208 samples).

\subsection{Prompt-Based Few-Shot Fine-Tuning}

Motivated by \citet{ullah2025prompt}, we frame hate speech detection as a cloze
task to leverage XLM-R's pre-trained MLM head rather than a randomly initialised
classification head.

\textbf{Template.} Each tweet is wrapped as:
\textit{``\{tweet\} This text is \texttt{[MASK]}.''}

\textbf{Verbalizer.}  Single-token XLM-R words are mapped to classes:
non-hate $\leftarrow$ \{normal, fine, okay, good\};
hate $\leftarrow$ \{hate, bad, wrong, vile, terrible\}.

\textbf{Loss.} At the \texttt{[MASK]} position we take log-sum-exp over each
class's verbalizer logits and apply cross-entropy.
Inference assigns the class whose verbalizer tokens receive higher total probability.

\textbf{Setup.} We use the XLM-R-large English checkpoint as the starting point,
train for 10 epochs with LR~$2\times10^{-5}$, batch size~8, and no class weighting.
Experiments use $K\in\{8,16,32,64\}$ stratified examples per class, repeated with
3 random seeds; we report mean~$\pm$~std.

\subsection{Transfer Protocols}

\textbf{Zero-shot:} English-trained classifier applied directly to Roman Urdu.

\textbf{Standard few-shot:} XLM-R-large English checkpoint fine-tuned on $K$-shot
Urdu with a classification head, LR~$2\times10^{-5}$, 5~epochs, balanced weights
$[1.0, 1.0]$ (re-calibrated for the Urdu distribution).

\textbf{Supervised upper bound:} XLM-R-large fine-tuned on the full Urdu training
set (7{,}208 samples), LR~$1\times10^{-5}$, 5~epochs, balanced weights.

\subsection{Evaluation Metrics}

We report hate~F1 (primary), macro~F1, non-hate~F1, and false positive rate
(FPR~$= \mathrm{FP}/(\mathrm{FP}+\mathrm{TN})$).

\subsection{Fairness Analysis}

We compute group-conditional FPR for text mentioning gender, religion, and
ethnicity terms via keyword lexicons.
For group~$G$: $\mathrm{FPR}_G = \mathrm{FP}_G / (\mathrm{FP}_G + \mathrm{TN}_G)$.

% -----------------------------------------------------------------------
\section{Results}
\label{sec:results}

\subsection{English Dev Set Performance}

Table~\ref{tab:en_results} shows English results.
XLM-R-large achieves the best hate~F1 (0.509) and macro~F1 (0.740), confirming
that model capacity improves cross-lingual representations.
The class-weighted loss successfully prevents collapse to all-non-hate on the
severely imbalanced English training set.

\begin{table}[h]
\centering
\caption{English dev set results (best epoch by hate~F1).}
\label{tab:en_results}
\vskip 0.05in
\begin{small}
\begin{tabular}{lcccc}
\toprule
Model & Hate F1 & NH F1 & Mac F1 & FPR \\
\midrule
mBERT          & 0.461 & 0.962 & 0.711 & 0.051 \\
XLM-R-base     & 0.494 & 0.964 & 0.729 & 0.049 \\
XLM-R-large    & \textbf{0.509} & \textbf{0.971} & \textbf{0.740} & \textbf{0.030} \\
\bottomrule
\end{tabular}
\end{small}
\end{table}

\subsection{Zero-Shot Transfer to Roman Urdu}

Table~\ref{tab:zs_results} shows zero-shot results.
All models achieve hate~F1~$<$~0.07, below the majority-class baseline macro~F1
of 0.356.
The failure mode is low hate recall ($<$~0.03): models predict nearly all samples
as non-hate, reflecting the strong non-hate bias learned from the imbalanced English
training set.
Notably, XLM-R-large achieves the \emph{lowest} zero-shot hate~F1 (0.005) ---
the larger model has a stronger EN-calibrated bias, making it even more conservative
on unseen Urdu text (RQ1, RQ2).

\begin{table}[h]
\centering
\caption{Zero-shot transfer to Roman Urdu dev set (800 samples).}
\label{tab:zs_results}
\vskip 0.05in
\begin{small}
\begin{tabular}{lcccc}
\toprule
Model & Hate F1 & NH F1 & Mac F1 & FPR \\
\midrule
Majority baseline  & 0.000 & 0.713 & 0.356 & 0.000 \\
mBERT              & 0.061 & 0.626 & 0.344 & 0.038 \\
XLM-R-base         & 0.058 & 0.630 & 0.344 & 0.027 \\
XLM-R-large        & 0.005 & 0.634 & 0.319 & 0.003 \\
\bottomrule
\end{tabular}
\end{small}
\end{table}

\subsection{Standard Few-Shot Transfer}

Table~\ref{tab:fs_results} shows few-shot results for XLM-R-large.
All runs converge to hate~F1~$\approx$~0.697, a degenerate solution where the
model predicts almost all samples as hate (FPR~$\approx$~1.0, non-hate~F1~$\approx$~0).
Even with balanced class weights $[1.0, 1.0]$ recalibrated for Urdu's
53.5\% hate prevalence, the collapse persists identically for $K \le 32$.
This is \textbf{not} a calibration artefact --- it is a fundamental limitation of
standard classification fine-tuning when $K \le 32$ Urdu examples cannot redirect
XLM-R-large's English-calibrated decision boundary (RQ3).

\begin{table}[h]
\centering
\caption{Standard few-shot results (XLM-R-large, mean $\pm$ std, 3 seeds).}
\label{tab:fs_results}
\vskip 0.05in
\begin{small}
\begin{tabular}{lccc}
\toprule
$K$ & Hate F1 & Mac F1 & FPR \\
\midrule
0 (zero-shot) & 0.005          & 0.319          & 0.003 \\
8             & 0.697 $\pm$ 0.000 & 0.350 $\pm$ 0.001 & $\approx$1.0 \\
16            & 0.697 $\pm$ 0.000 & 0.351 $\pm$ 0.001 & $\approx$1.0 \\
32            & 0.697 $\pm$ 0.000 & 0.349 $\pm$ 0.001 & $\approx$1.0 \\
64            & 0.715 $\pm$ 0.009 & 0.528 $\pm$ 0.091 & 0.777 \\
\bottomrule
\end{tabular}
\end{small}
\end{table}

\subsection{Prompt-Based Few-Shot Fine-Tuning}

Table~\ref{tab:prompt_results} shows prompt-based results.
The MLM verbalizer approach avoids the random-initialisation problem of the
classification head: even at $K=8$, macro~F1 is 0.474 vs.\ 0.350 for standard
fine-tuning.
At $K=32$ prompt-based fine-tuning achieves macro~F1~$=$~0.634 and
FPR~$=$~0.550, vs.\ macro~F1~$=$~0.349 and FPR~$\approx$~1.0 for standard
fine-tuning.
At $K=64$ it reaches hate~F1~$=$~0.787 and macro~F1~$=$~0.745 (RQ3).

\begin{table}[h]
\centering
\caption{Prompt-based few-shot results (XLM-R-large, mean $\pm$ std, 3 seeds).}
\label{tab:prompt_results}
\vskip 0.05in
\begin{small}
\begin{tabular}{lccc}
\toprule
$K$ & Hate F1 & Mac F1 & FPR \\
\midrule
8  & 0.690 $\pm$ 0.021 & 0.474 $\pm$ 0.121 & 0.801 \\
16 & 0.697 $\pm$ 0.007 & 0.496 $\pm$ 0.054 & 0.803 \\
32 & 0.737 $\pm$ 0.026 & 0.634 $\pm$ 0.099 & 0.550 \\
64 & \textbf{0.787 $\pm$ 0.012} & \textbf{0.745 $\pm$ 0.038} & \textbf{0.355} \\
\bottomrule
\end{tabular}
\end{small}
\end{table}

\subsection{Classical ML Baselines and Supervised Upper Bound}

Table~\ref{tab:classical_results} shows results for classical ML classifiers
trained on the \emph{full} Urdu training set (7{,}208 samples) and the supervised
XLM-R-large upper bound.
TF-IDF$+$Fuzzy$+$LR achieves hate~F1~$=$~0.857, outperforming all few-shot
neural approaches.
This confirms that target-language lexical coverage matters more than cross-lingual
representations when labelled Urdu data is available.
Full supervised fine-tuning of XLM-R-large further improves to hate~F1~$=$~0.895
and macro~F1~$=$~0.885, quantifying the remaining gap that few-shot and
prompt-based methods must close.

\begin{table}[h]
\centering
\caption{Classical ML and supervised results on Roman Urdu dev set.}
\label{tab:classical_results}
\vskip 0.05in
\begin{small}
\begin{tabular}{lcccc}
\toprule
Model & Hate F1 & NH F1 & Mac F1 & FPR \\
\midrule
TF-IDF + LR           & 0.854 & 0.817 & 0.835 & 0.220 \\
TF-IDF + RF           & 0.851 & 0.788 & 0.819 & 0.301 \\
TF-IDF+Fuzzy+LR       & \textbf{0.857} & 0.824 & \textbf{0.841} & \textbf{0.207} \\
TF-IDF+Fuzzy+RF       & 0.837 & 0.779 & 0.808 & 0.290 \\
\midrule
XLM-R-large (full UR) & \textbf{0.895} & \textbf{0.876} & \textbf{0.885} & \textbf{0.140} \\
\bottomrule
\end{tabular}
\end{small}
\end{table}

\subsection{Fairness Analysis}

Table~\ref{tab:fairness} shows group-conditional FPR for XLM-R-base zero-shot
predictions on Roman Urdu (XLM-R-large zero-shot was too conservative to produce
meaningful FPR comparisons, predicting almost exclusively non-hate).

\begin{table}[h]
\centering
\caption{Group-conditional FPR (XLM-R-base, zero-shot Urdu).}
\label{tab:fairness}
\vskip 0.05in
\begin{small}
\begin{tabular}{lrcc}
\toprule
Group & $n$ & FPR & Hate F1 \\
\midrule
Gender    &  33 & 0.059 & 0.000 \\
Religion  &  30 & \textbf{0.143} & 0.615 \\
Ethnicity &   4 & 0.000 & 0.000 \\
Overall   & 800 & 0.027 & 0.058 \\
\bottomrule
\end{tabular}
\end{small}
\end{table}

Religion-mentioning text shows a 5$\times$ higher FPR than overall (0.143
vs.\ 0.027), suggesting the model falsely flags non-hate religious content as
hate at a significantly elevated rate.
This likely reflects patterns in the English training data where religious
terminology co-occurs with hate speech (RQ4).

% -----------------------------------------------------------------------
\section{Ablation Studies}
\label{sec:ablation}

\paragraph{A1: Model backbone.}
Tables~\ref{tab:en_results} and~\ref{tab:zs_results} reveal a clear capacity
hierarchy on English (mBERT~$<$~XLM-R-base~$<$~XLM-R-large in macro~F1 and hate~F1),
but this ordering \emph{reverses} on zero-shot Urdu: XLM-R-large achieves the
lowest hate~F1 (0.005) because stronger English fine-tuning produces a more rigid,
EN-calibrated decision boundary.
Larger models are harder to redirect toward an unseen target language with only
zero examples.

\paragraph{A2: Training method vs.\ $K$-shot data.}
Table~\ref{tab:ablation} compares standard fine-tuning and prompt-based fine-tuning
at matched $K$, alongside the fully-supervised baselines.
Prompt-based fine-tuning strictly dominates standard fine-tuning at every $K$
in macro~F1 and FPR.
The key advantage is that the pre-trained MLM head provides meaningful token
probability estimates from the first iteration, whereas a randomly initialised
classification head must first be learned from scratch on $\le 64$ examples.
Nevertheless, both few-shot approaches fall substantially short of the classical
ML baselines trained on all 7{,}208 Urdu samples, confirming that target-language
lexical supervision remains the binding constraint.

\begin{table}[h]
\centering
\caption{Method comparison at matched conditions (XLM-R-large).}
\label{tab:ablation}
\vskip 0.05in
\begin{small}
\begin{tabular}{llccc}
\toprule
Method & $K$ & Hate F1 & Mac F1 & FPR \\
\midrule
Standard FT    &  8  & 0.697 & 0.350 & 0.998 \\
Prompt-based   &  8  & 0.690 & 0.474 & 0.801 \\
\midrule
Standard FT    & 32  & 0.697 & 0.349 & 1.000 \\
Prompt-based   & 32  & 0.737 & 0.634 & 0.550 \\
\midrule
Standard FT    & 64  & 0.715 & 0.528 & 0.777 \\
Prompt-based   & 64  & 0.787 & 0.745 & 0.355 \\
\midrule
TF-IDF+Fuzzy+LR & all & 0.857 & 0.841 & 0.207 \\
XLM-R-l (sup.) & all & \textbf{0.895} & \textbf{0.885} & \textbf{0.140} \\
\bottomrule
\end{tabular}
\end{small}
\end{table}

\paragraph{A3: Class weighting in few-shot.}
We originally used weights $[1.0, 3.0]$ tuned for the English distribution (5.2\%
hate), then re-ran with balanced $[1.0, 1.0]$ weights matching Urdu's 53.5\%
hate prevalence.
The all-hate collapse persisted identically at $K \le 32$ under both settings,
ruling out calibration mismatch as the root cause and pointing to the fundamental
difficulty of redirecting a large model's decision boundary from only a handful
of target-language examples.

% -----------------------------------------------------------------------
\section{Discussion}
\label{sec:discussion}

\paragraph{Annotation disagreement.}
The Davidson et al.\ (2017) dataset has documented inter-annotator disagreement,
particularly for offensive but non-hateful content.
The 5.2\% hate label rate may underrepresent actual hate speech if annotators
systematically coded borderline cases as ``offensive'' rather than ``hate.''
This labelling conservatism propagates into the model and contributes to the
near-zero recall observed on the balanced Roman Urdu corpus.

\paragraph{Cross-lingual mismatch and path forward.}
Roman Urdu is not well-represented in XLM-R's Common Crawl pre-training data
(which indexes standard Urdu in Nastaliq script).
Our results quantify the resulting gap: even with prompt-based fine-tuning on
64 Urdu examples, hate~F1 reaches 0.787 vs.\ 0.895 for full supervision ---
a 0.108 gap that represents the cost of cross-lingual transfer with minimal
target-language data.
The strong performance of TF-IDF$+$FuzzyWuzzy (hate~F1~$=$~0.857) on full Urdu
data shows that lexical approaches remain highly competitive for morphologically
rich, domain-specific text when sufficient labelled target-language data is
available.

% -----------------------------------------------------------------------
\section{Conclusion}
\label{sec:conclusion}

We conducted a comprehensive study of cross-lingual hate speech transfer from
English to Roman Urdu using multilingual transformers.
Our key findings are:
(1)~Zero-shot transfer fails for all models (hate~F1~$<$~0.07), with larger
models exhibiting a stronger EN-calibrated bias.
(2)~Standard few-shot fine-tuning collapses to all-hate predictions for
$K \le 32$, a failure not attributable to class-weight miscalibration.
(3)~Prompt-based fine-tuning with an MLM verbalizer substantially mitigates
collapse, improving macro~F1 by $+$0.396 over standard fine-tuning at $K=32$.
(4)~Classical TF-IDF$+$FuzzyWuzzy baselines and supervised XLM-R-large set
hate~F1 upper bounds of 0.857 and 0.895 respectively.
(5)~Fairness analysis reveals 5$\times$ elevated FPR on religion-mentioning text.

These results highlight that cross-lingual transfer to low-resource, informal
scripts remains an open challenge.
Future work should explore data augmentation, cross-lingual few-shot alignment,
and larger-scale Roman Urdu pre-training to close the remaining gap.

% -----------------------------------------------------------------------
\section*{AI Tools Statement}
Claude Sonnet 4.6 (Anthropic) assisted with code generation, debugging, and
report drafting. All experimental results are produced by our code running on
our hardware (NVIDIA RTX 4070 GPU). All claims and results were verified by the
authors.

% -----------------------------------------------------------------------
\bibliography{references}
\bibliographystyle{icml2025}

\end{document}
